\chapter{👨‍👩‍👧‍👦 家族 — La Familia}
\unidadheader{10}{家族}{La Familia}{Hablar de tu familia}

\begin{tcolorbox}[vocabbox, title={📌 Vocabulario Nuevo — うち (mi familia) vs そと (familia ajena)}]
\begin{tabularx}{\linewidth}{lll}
\toprule
\textbf{Familiar} & \textbf{Mi familia (うち)} & \textbf{Familia ajena (そと)} \\
\midrule
padre & \ruby{父}{ちち} & お\ruby{父}{とう}さん \\
madre & \ruby{母}{はは} & お\ruby{母}{かあ}さん \\
hermano mayor & \ruby{兄}{あに} & お\ruby{兄}{にい}さん \\
hermana mayor & \ruby{姉}{あね} & お\ruby{姉}{ねえ}さん \\
hermano menor & \ruby{弟}{おとうと} & \ruby{弟}{おとうと}さん \\
hermana menor & \ruby{妹}{いもうと} & \ruby{妹}{いもうと}さん \\
abuelo & \ruby{祖父}{そふ} & お\ruby{祖父}{じい}さん \\
abuela & \ruby{祖母}{そぼ} & お\ruby{祖母}{ばあ}さん \\
\bottomrule
\end{tabularx}
\end{tcolorbox}

\subsection*{🗣️ Frases Clave}

\frase{\ruby{家族}{かぞく}は\ruby{何人}{なんにん}ですか？}{¿Cuántos miembros tiene tu familia?}
\frase{\ruby{四人家族}{よにんかぞく}です。}{Somos familia de 4.}
\frase{\ruby{兄}{あに}が\ruby{一人}{ひとり}います。}{Tengo un hermano mayor.}
\frase{\ruby{兄弟}{きょうだい}はいますか？}{¿Tienes hermanos?}
\frase{一\ruby{人}{り}っ\ruby{子}{こ}です。}{Soy hijo/a único/a.}

\subsection*{⚙️ Gramática}

\patron{Cuántos hay: contador de personas}
  {〜人 (にん) / ひとり・ふたり}
  {1人 ひとり, 2人 ふたり, 3人以上 〜にん
  \ej{\ruby{家族}{かぞく}は\ruby{五人}{ごにん}います。}
    {En mi familia somos 5 personas.}
  \ej{\ruby{弟}{おとうと}が\ruby{二人}{ふたり}います。}
    {Tengo dos hermanos menores.}}

\patron{Presentar a la familia (うち vs そと)}
  {}
  {Hablar de MI familia → forma humilde (うち): 父、母、兄…\\
  Hablar de la familia AJENA → forma honorífica (そと): お父さん、お母さん…
  \ej{これは\ruby{父}{ちち}です。}{Este es mi padre. (hablando de propio padre)}
  \ej{お\ruby{父}{とう}さんはお\ruby{元気}{げんき}ですか？}
    {¿Está bien su padre? (preguntando por el de otro)}}

\begin{tcolorbox}[ejerciciobox, title={📝 Ejercicios Rápidos}]
\begin{enumerate}
  \item Di cómo se llama tu familia usando うち vocabulary.
  \item Pregunta a alguien cuántos hermanos tiene.
  \item ¿Cómo dices "Tengo una hermana mayor y un hermano menor"?
\end{enumerate}
\vspace{0.4em}
\textbf{Respuestas:}
1) (personal)\quad
2) \ruby{兄弟}{きょうだい}は\ruby{何人}{なんにん}いますか？\quad
3) \ruby{姉}{あね}が\ruby{一人}{ひとり}と\ruby{弟}{おとうと}が\ruby{一人}{ひとり}います。
\end{tcolorbox}

\begin{tcolorbox}[kanjibox, title={🀄 Kanjis de la Unidad}]
\begin{tabularx}{\linewidth}{cllXX}
\toprule
\textbf{Kanji} & \textbf{On} & \textbf{Kun} & \textbf{Significado} & \textbf{Vocab} \\
\midrule
\kanjirow{父}{ふ}{ちち}{padre}{\ruby{父}{ちち} / お\ruby{父}{とう}さん}
\kanjirow{母}{ぼ}{はは}{madre}{\ruby{母}{はは} / お\ruby{母}{かあ}さん}
\kanjirow{兄}{けい}{あに}{hermano mayor}{\ruby{兄}{あに} / お\ruby{兄}{にい}さん}
\kanjirow{姉}{し}{あね}{hermana mayor}{\ruby{姉}{あね} / お\ruby{姉}{ねえ}さん}
\kanjirow{弟}{てい}{おとうと}{hermano menor}{\ruby{弟}{おとうと}}
\kanjirow{妹}{まい}{いもうと}{hermana menor}{\ruby{妹}{いもうと}}
\bottomrule
\end{tabularx}
\end{tcolorbox}
